% This document was generated by an LLM.

\documentclass{essay}

\title{An Anatomy of Precision:\\The Epsilon--Delta Rigor of Limits and Continuity}
\subtitle{A Mathematical Essay on Quantifiers, Inequalities, and Edge Cases}
\author{}
\date{}

\begin{document}

\maketitle

\begin{abstract}
\noindent
The Cauchy--Weierstrass $(\epsilon, \delta)$-definitions of function limits and continuity are among the basic tools of analysis. While introductory calculus presents these definitions as standard technical formalisms, their small details matter: the quantifiers, inequalities, and exclusions are not interchangeable. This essay examines the definitions of function limits and continuity, with special attention to why $0 < |x - a|$ is required for limits but absent in continuity, why strict inequalities are chosen over non-strict inequalities in quantifiers, and what goes wrong if these conditions are modified or relaxed.
\end{abstract}

\tableofcontents

\section{Introduction}

In early calculus, a limit $\lim_{x \to a} f(x) = L$ is often intuitively described as: \textit{``as $x$ gets arbitrarily close to $a$, $f(x)$ gets arbitrarily close to $L$.''} While visually appealing, this heuristic relies on kinematic concepts (``gets close to'', ``approaches'') that invite logical paradoxes and ambiguity.

During the $19^\text{th}$ century, Augustin-Louis Cauchy and Karl Weierstrass arithmeticized analysis by replacing dynamic motion with static logical quantifiers. The modern definitions convert topological proximity into precise real-number inequalities. However, it is easy to take the structural details of these definitions for granted. Why do we write $0 < |x - a| < \delta$ instead of $|x - a| < \delta$? Why do we require $\epsilon > 0$ strictly? What breaks if we replace $< \epsilon$ with $\le \epsilon$?

This essay analyzes these subtleties and shows that some small-looking changes alter the meaning of the definitions.

\section{A Prototype: Equality by Arbitrary Precision}

Before looking at limits, it is useful to isolate a simpler logical pattern that already contains the heart of $(\epsilon,\delta)$ reasoning. The pattern is this: if a nonnegative quantity can be made smaller than every positive tolerance, then that quantity must be zero.

\begin{theorem}[Equality by Arbitrary Precision]
\label{thm:equality-by-arbitrary-precision}
For any real numbers $a,b \in \mathbb{R}$,
\[
a=b
\quad\Longleftrightarrow\quad
\forall \epsilon>0,\ |a-b|<\epsilon.
\]
\end{theorem}

\begin{proof}
We prove both directions.

\textbf{Direction 1 ($a=b \implies \forall \epsilon>0,\ |a-b|<\epsilon$):}
Assume $a=b$. Then $a-b=0$, so $|a-b|=0$. Let $\epsilon>0$ be arbitrary. Since $0<\epsilon$, we have $|a-b|=0<\epsilon$. Therefore
\[
\forall \epsilon>0,\ |a-b|<\epsilon.
\]

\textbf{Direction 2 ($\forall \epsilon>0,\ |a-b|<\epsilon \implies a=b$):}
Assume
\[
\forall \epsilon>0,\ |a-b|<\epsilon.
\]
We will show that $a=b$. Suppose, for contradiction, that $a\neq b$. Then $|a-b|>0$. Choose
\[
\epsilon_0=\frac{|a-b|}{2}.
\]
Since $|a-b|>0$, we have $\epsilon_0>0$. By the assumed statement, applied to this particular positive tolerance $\epsilon_0$, we get
\[
|a-b|<\epsilon_0=\frac{|a-b|}{2}.
\]
But this is impossible for a positive number $|a-b|$: it would imply
\[
\frac{|a-b|}{2}<0.
\]
This contradiction shows that our supposition $a\neq b$ was false. Hence $a=b$.
\end{proof}

This theorem explains why the phrase ``for every $\epsilon>0$'' is so powerful. It does not merely say that $|a-b|$ is small. It says that no positive tolerance, however tiny, can still detect any separation between $a$ and $b$. The only nonnegative distance with that property is zero.

The $(\epsilon,\delta)$-definition of a limit adds one extra layer to this prototype. Instead of directly proving that a fixed distance $|a-b|$ is smaller than every $\epsilon$, it says:
\[
\forall \epsilon>0,\ \exists \delta>0
\]
such that whenever $x$ is within the permitted $\delta$-neighborhood of $a$, the output distance $|f(x)-L|$ is smaller than that chosen $\epsilon$. In other words, the $\delta$ clause is a mechanism for forcing the output error below an arbitrary tolerance.

Thus a limit proof usually has the same skeleton as Theorem~\ref{thm:equality-by-arbitrary-precision}, but with a dependency inserted:
\[
\text{given an arbitrary }\epsilon>0,
\quad
\text{choose a suitable }\delta>0,
\quad
\text{then prove } |f(x)-L|<\epsilon
\]
for every $x$ satisfying the domain and closeness conditions. The theorem above is the degenerate case where there is no input variable and no $\delta$ to choose: the relevant distance is already fixed.

\section{The Formal Definitions}

Let $D \subseteq \mathbb{R}$ be the domain of a real-valued function $f: D \to \mathbb{R}$, and let $a \in \mathbb{R}$ be a limit point (or accumulation point) of $D$.

\begin{definition}[Limit of a Function]
\label{def:limit}
We say that the \textbf{limit of $f(x)$ as $x$ approaches $a$ is $L$}, denoted $\lim_{x \to a} f(x) = L$, if:
\[
\forall \epsilon > 0, \, \exists \delta > 0, \, \forall x \in D, \, \left( 0 < |x - a| < \delta \implies |f(x) - L| < \epsilon \right).
\]
\end{definition}

Now consider the definition of continuity at a point. Let $a \in D$.

\begin{definition}[Continuity at a Point]
\label{def:continuity}
We say that $f$ is \textbf{continuous at $a$} if:
\[
\forall \epsilon > 0, \, \exists \delta > 0, \, \forall x \in D, \, \left( |x - a| < \delta \implies |f(x) - f(a)| < \epsilon \right).
\]
\end{definition}

Comparing Definition~\ref{def:limit} and Definition~\ref{def:continuity}, two structural differences immediately emerge:
\begin{enumerate}
    \item \textbf{The explicit value of $L$ vs. $f(a)$}: In the limit definition, $L$ is a candidate value not necessarily tied to $f(a)$ (indeed, $f(a)$ need not even be defined). In continuity, $f(a)$ is strictly required to exist, and $L = f(a)$.
    \item \textbf{The lower bound $0 < |x - a|$}: The limit definition explicitly excludes $x = a$, whereas the definition of continuity includes $x = a$.
\end{enumerate}

\section{Excluding Zero: The Role of $0 < |x - a|$}

The condition $0 < |x - a|$ is equivalent to stating $x \neq a$. Its inclusion in the definition of a limit, and its omission in the definition of continuity, is not accidental.

\subsection{Why Exclude $x = a$ in Limits?}

The limit $\lim_{x \to a} f(x)$ seeks to describe the \emph{local behavior of $f$ near $a$}, completely independent of the \emph{value of $f$ at $a$}.

Suppose we removed the constraint $0 < |x - a|$ and defined a "modified limit" $L'$ using:
\[
|x - a| < \delta \implies |f(x) - L'| < \epsilon.
\]
If $x = a$ were allowed, then $|a - a| = 0 < \delta$ holds for every $\delta > 0$. The implication would then require:
\[
|f(a) - L'| < \epsilon \quad \forall \epsilon > 0.
\]
Since the only non-negative real number less than every positive $\epsilon$ is $0$, this forces $|f(a) - L'| = 0$, meaning $L'$ \emph{must} equal $f(a)$.

This hypothetical alteration would immediately break standard examples in calculus:
\begin{itemize}
    \item \textbf{Functions with removable discontinuities would have no limits.} Consider $f(x) = \frac{x^2 - 1}{x - 1}$ for $x \neq 1$, with $f(1) = 5$. Intuitively and standardly, $\lim_{x \to 1} f(x) = 2$. However, under the modified definition, no limit could exist because $f(1) = 5 \neq 2$.
    \item \textbf{Functions undefined at $a$ could not have limits.} The standard derivative definition is:
    \[
    f'(a) = \lim_{h \to 0} \frac{f(a+h) - f(a)}{h}.
    \]
    The difference quotient $g(h) = \frac{f(a+h) - f(a)}{h}$ is explicitly undefined at $h = 0$. If $h = 0$ were required in the limit quantifier, no derivative could ever be evaluated!
\end{itemize}

Thus, excluding $x = a$ decouples local trend from point behavior.

\subsection{Why Include $x = a$ in Continuity?}

For continuity, the situation is reversed. Continuity asserts that the local trend \emph{matches} the point value.

If we evaluate the continuity condition at $x = a$:
\[
|a - a| < \delta \implies |f(a) - f(a)| < \epsilon \iff 0 < \delta \implies 0 < \epsilon.
\]
For $x = a$, the statement $|f(a) - f(a)| < \epsilon$ reduces to $0 < \epsilon$, which is vacuously true for all $\epsilon > 0$. Thus, explicitly writing $0 < |x - a| < \delta$ or $|x - a| < \delta$ in the definition of continuity produces logically equivalent statements. Including $x = a$ simply simplifies the notation and emphasizes that $f(a)$ must be defined.

\section{Strict vs. Non-Strict Inequalities in Quantifiers}

A subtle point of confusion for students is why we use strict inequalities ($\epsilon > 0$, $\delta > 0$, $|x - a| < \delta$, $|f(x) - L| < \epsilon$) rather than weak inequalities ($\ge, \le$). Does changing strict inequalities to weak ones change the mathematical meaning?

Let us test all permutations systematically.

\subsection{Can we replace $|f(x) - L| < \epsilon$ with $|f(x) - L| \le \epsilon$?}

\begin{theorem}
The definition of limit with $|f(x) - L| \le \epsilon$ is equivalent to Definition~\ref{def:limit}.
\end{theorem}

\begin{proof}
Let Definition A be the standard definition using $< \epsilon$, and Definition B be the modified definition using $\le \epsilon$.

\textbf{Direction 1 (A $\implies$ B):} Assume $\lim_{x \to a} f(x) = L$ under Definition A. Given $\epsilon > 0$, by Definition A there exists $\delta > 0$ such that $0 < |x - a| < \delta \implies |f(x) - L| < \epsilon$. Since $|f(x) - L| < \epsilon$ implies $|f(x) - L| \le \epsilon$, Definition B holds with the same $\delta$.

\textbf{Direction 2 (B $\implies$ A):} Assume Definition B holds. Given $\epsilon > 0$, choose $\epsilon' = \frac{\epsilon}{2} > 0$. By Definition B, there exists $\delta > 0$ such that $0 < |x - a| < \delta \implies |f(x) - L| \le \epsilon' = \frac{\epsilon}{2}$. Since $\frac{\epsilon}{2} < \epsilon$, we have $|f(x) - L| < \epsilon$. Thus Definition A holds.
\end{proof}

Hence, using $\le \epsilon$ instead of $< \epsilon$ changes nothing mathematically. The strict inequality $< \epsilon$ is preferred purely for aesthetic symmetry with $|x - a| < \delta$.

\subsection{Can we replace $0 < |x - a| < \delta$ with $0 < |x - a| \le \delta$?}

\begin{theorem}
The definition of limit with $0 < |x - a| \le \delta$ is equivalent to Definition~\ref{def:limit}.
\end{theorem}

\begin{proof}
If a $\delta_1 > 0$ satisfies the strict inequality condition, taking $\delta_2 = \frac{\delta_1}{2} > 0$ satisfies the weak inequality condition because $0 < |x - a| \le \delta_2 \implies 0 < |x - a| < \delta_1$. Conversely, if $\delta_1$ works for weak inequality, it automatically works for strict inequality since $0 < |x - a| < \delta_1 \implies 0 < |x - a| \le \delta_1$.
\end{proof}

\subsection{What if we allow $\epsilon = 0$ ($\epsilon \ge 0$)?}

What happens if we modify the quantifier from $\forall \epsilon > 0$ to $\forall \epsilon \ge 0$?

Suppose we require:
\[
\forall \epsilon \ge 0, \, \exists \delta > 0, \, 0 < |x - a| < \delta \implies |f(x) - L| \le \epsilon.
\]
Setting $\epsilon = 0$, the condition demands that there exists $\delta > 0$ such that for all $x$ in the punctured neighborhood $0 < |x - a| < \delta$:
\[
|f(x) - L| \le 0 \implies f(x) = L.
\]
This forces $f(x)$ to be \emph{strictly constant} in an open punctured neighborhood around $a$.

Thus, allowing $\epsilon = 0$ makes the definition far too strong. Non-constant functions like $f(x) = x$ would fail to have a limit anywhere! The strict positivity $\epsilon > 0$ is what allows $f(x)$ to approximate $L$ without forcing immediate equality.

\subsection{What if we allow $\delta = 0$ ($\delta \ge 0$)?}

Suppose the definition allowed choosing $\delta = 0$.
If $\delta = 0$, the condition $0 < |x - a| < 0$ is a contradiction—no real number $x$ can satisfy it. The antecedent of the implication:
\[
0 < |x - a| < 0 \implies |f(x) - L| < \epsilon
\]
becomes \emph{vacuously true} for \emph{every} real number $L$ and every function $f$.

If $\delta = 0$ were allowed, every function would have \emph{every} real number as its limit at every point, so limits would no longer be unique. Hence, $\delta$ must be strictly positive ($\delta > 0$).

\section{Summary Matrix of Variant Definitions}

To synthesize these logical dependencies, the table below summarizes the effect of altering individual components of the $(\epsilon, \delta)$-definition for $\lim_{x \to a} f(x) = L$:

\begin{table}[h!]
\centering
\small
\begin{tabular}{|l|l|l|}
\hline
\textbf{Modification} & \textbf{Mathematical Consequence} & \textbf{Status} \\ \hline
Replace $< \epsilon$ with $\le \epsilon$ & Equivalent. & Valid alternative. \\ \hline
Replace $< \delta$ with $\le \delta$ & Equivalent. & Valid alternative. \\ \hline
Remove $0 < |x - a|$ & Forces $L = f(a)$. Breaks derivatives. & Invalid. \\ \hline
Allow $\epsilon \ge 0$ & Forces $f(x) = L$ locally (constant function). & Invalid. \\ \hline
Allow $\delta \ge 0$ & Vacuously true for all $L$. Limits not unique. & Invalid. \\ \hline
Swap quantifiers: $\exists \delta \forall \epsilon$ & Forces $f(x) = L$ identically near $a$. & Invalid. \\ \hline
\end{tabular}
\caption{Impact of small formal changes on the limit definition.}
\end{table}

\section{Conclusion}

The $(\epsilon, \delta)$-definition of a limit is carefully tuned. The inclusion of $0 < |x - a|$ separates the limit from the function's value at the point, which is what makes removable discontinuities and derivative quotients possible. The strict positivity of $\epsilon$ and $\delta$ leaves room for approximation without collapsing the definition into exact equality or vacuous truth. Choices that can look like minor notational preferences are often doing real logical work.

\end{document}
